\par\Needspace{6\baselineskip}\medskip
\hypertarget{native-211-article156}{}\pdfbookmark[2]{156.}{article-156}
\begin{center}156.\end{center}
We shall take our \textit{second} example from Pallas. We excerpt the following observations made in Milan from the scholarly correspondence of the distinguished \textsc{de Zach}, Vol.\ XIV. p.\ 90.
\begin{center}
\begingroup\footnotesize\setlength{\tabcolsep}{4pt}\renewcommand{\arraystretch}{1.2}
\begin{tabular}{|l|r|l|}\hline
Milan mean time & App.\ right asc. & App.\ decl.\\\hline
1805 Nov.\ 5.\ $14^{\mathrm h}\,14^{\mathrm m}\,4^{\mathrm s}$ & $78^{0}\,20'\,37''{,}8$ & $27^{0}\,16'\,56''{,}7$ S.\\
\hspace{2.4em}Dec.\ 6.\ $11\quad51\quad27$ & $73\quad8\quad48{,}8$ & $32\quad52\quad44{,}3$\\
1806 Jan.\ 15.\ $8\quad50\quad36$ & $67\quad14\quad11{,}1$ & $28\quad38\quad8{,}1$\\\hline
\end{tabular}\endgroup
\end{center}

Here we shall take the equator as the fundamental plane in place of the ecliptic and carry out the computation as though the orbit were still wholly unknown. First we take the following from the solar tables for the given times:
\begin{center}
\begingroup\footnotesize\setlength{\tabcolsep}{4pt}\renewcommand{\arraystretch}{1.2}
\begin{tabular}{|l|r|r|r|}\hline
& \shortstack{Sun's longitude from\\mean equinox} & Distance from Earth & \shortstack{Sun's\\latitude}\\\hline
Nov.\ 5 & $223^{0}\,14'\,7''{,}61$ & $0{,}990\,4311$ & $+0''{,}59$\\
Dec.\ 6 & $254\quad28\quad42{,}59$ & $0{,}984\,6753$ & $+0{,}12$\\
Jan.\ 15 & $295\quad5\quad47{,}62$ & $0{,}983\,8153$ & $-0{,}19$\\\hline
\end{tabular}\endgroup
\end{center}

By adding the precessions $+7''{,}59$, $+3''{,}36$, $-2''{,}11$, we reduce the longitudes of the Sun to the beginning of 1806. Then, using the mean obliquity $23^{0}\,27'\,53''{,}53$ and taking due account of the latitudes, we derive the right ascensions and declinations. We thus find
\SourceLeaf{212}
\begin{center}
\begingroup\small\setlength{\tabcolsep}{5pt}\renewcommand{\arraystretch}{1.2}
\begin{tabular}{|l|r|l|}\hline
& Sun's right ascension & Sun's declination\\\hline
Nov.\ 5 & $220^{0}\,46'\,44''{,}65$ & $15^{0}\,49'\,43''{,}94$ S.\\
Dec.\ 6 & $253\quad9\quad23{,}26$ & $22\quad33\quad39{,}45$\\
Jan.\ 15 & $297\quad2\quad51{,}11$ & $21\quad8\quad12{,}98$\\\hline
\end{tabular}\endgroup
\end{center}

These positions are referred to the center of the Earth and must therefore be reduced to the place of observation by adding the parallax, since the planet's positions cannot be freed from parallax. The right ascensions of the zenith to be used in this computation agree with the right ascensions of the planet (since the observations were made on the meridian itself), while the declination in every case is the altitude of the pole $=45^{0}\,28'$. This gives the following numbers:
\begin{center}
\begingroup\footnotesize\setlength{\tabcolsep}{3pt}\renewcommand{\arraystretch}{1.2}
\begin{tabular}{|l|r|l|r|}\hline
& Earth's right asc. & Earth's declination & Log.\ dist.\ from Sun\\\hline
Nov.\ 5 & $40^{0}\,46'\,48''{,}51$ & $15^{0}\,49'\,48''{,}59$ N. & $9{,}995\,8375$\\
Dec.\ 6 & $73\quad9\quad23{,}26$ & $22\quad33\quad42{,}83$ & $9{,}993\,3099$\\
Jan.\ 15 & $117\quad2\quad46{,}09$ & $21\quad8\quad17{,}29$ & $9{,}992\,9259$\\\hline
\end{tabular}\endgroup
\end{center}

The observed positions of Pallas must be freed from nutation and the aberration of the fixed stars, and then reduced to the beginning of 1806 by adding the precession. Under these headings the following corrections must be applied to the observed positions:
\begin{center}
\begingroup\footnotesize\setlength{\tabcolsep}{4pt}\renewcommand{\arraystretch}{1.2}
\begin{tabular}{|l|r|r||r|r||r|r|}\hline
& \multicolumn{2}{c||}{Observation I} & \multicolumn{2}{c||}{Observation II} & \multicolumn{2}{c|}{Observation III}\\\hline
& R.\ asc. & Decl. & R.\ asc. & Decl. & R.\ asc. & Decl.\\\hline
Nutation & $-12''{,}86$ & $-3''{,}08$ & $-12''{,}21$ & $-3''{,}42$ & $-13''{,}06$ & $-3''{,}75$\\
aberration & $-18{,}13$ & $-9{,}89$ & $-22{,}68$ & $-1{,}63$ & $-15{,}60$ & $+9{,}76$\\
precession & $+5{,}43$ & $+0{,}62$ & $+2{,}25$ & $+0{,}39$ & $-1{,}51$ & $-0{,}33$\\\hline
sum & $-25{,}56$ & $-12{,}35$ & $-32{,}64$ & $-4{,}66$ & $-30{,}17$ & $+5{,}68$\\\hline
\end{tabular}\endgroup
\end{center}

This gives the following positions of Pallas, on which the computation is to be based:
\begin{center}
\begingroup\small\setlength{\tabcolsep}{6pt}\renewcommand{\arraystretch}{1.2}
\begin{tabular}{|r|r|r|}\hline
Paris mean time & Right asc. & Declination\\\hline
Nov.\ $5{,}574\,074$ & $78^{0}\,20'\,12''{,}24$ & $-27^{0}\,17'\,9''{,}05$\\
$36{,}475\,035$ & $73\quad8\quad16{,}16$ & $-32\quad52\quad48{,}96$\\
$76{,}349\,444$ & $67\quad13\quad40{,}93$ & $-28\quad38\quad2{,}42$\\\hline
\end{tabular}\endgroup
\end{center}
